\documentclass[12pt]{article}

\usepackage[margin=1in]{geometry}
\usepackage{amsmath,amsthm,amssymb}
\usepackage{xeCJK,hyperref}
\usepackage{graphicx}
\usepackage{minted, algorithm, listings}
\usepackage{tikz}
\usepackage{tikz-feynman}
\usepackage{diagbox,fancybox}

% ==============================================================================
%  Commands
% ==============================================================================
\newcommand{\N}{\mathbb{N}}
\newcommand{\Z}{\mathbb{Z}}

% ==============================================================================
%  Title
% ==============================================================================
\title{Large Index Expansion for Feynman Integral Reduction: A Comprehensive Report}
\author{W \\
  Supervisor: M}
\date{\today}

\begin{document}

\maketitle

\begin{abstract}
  We present a comprehensive account of the \textbf{Large Index Expansion (LIE)} method
  for reducing Feynman integrals to a set of master integrals.
  The method expands the integrand in the limit $n\to\infty$ after the substitution
  $\nu_a\to n\theta_a+\nu_a$, reconstructs asymptotic expansion coefficients via
  integration-by-parts (IBP) relations, and builds reduction rules through a
  two-step procedure: (i) finding a minimal generating set, then (ii) constructing
  solvable block relations on a cone-shaped index set.
  We detail the complete theoretical framework, the block-recursive structure of
  expansion coefficients, a geometric classification of solution spaces, the
  asymptotic-solution completeness theorem grounded in microlocal analysis, and a
  high-performance C++ implementation using finite-field arithmetic that achieves a
  $20$--$30\times$ speedup over the prototype \emph{Mathematica} code.
\end{abstract}

% ==============================================================================
%  Section 1: Introduction
% ==============================================================================
\section{Introduction}\label{sec:intro}

In perturbative quantum field theory (pQFT), computing cross sections proceeds through
four stages:
\begin{enumerate}
  \item Generate Feynman diagrams from Feynman rules;
  \item Reduce tensor integrals to scalar Feynman integrals (FIs);
  \item Express each FI as a linear combination of a finite set of \textbf{master integrals} (MIs);
        \label{item:reduction}
  \item Evaluate the MIs.
\end{enumerate}
Step~\ref{item:reduction} --- \emph{integral reduction} --- is the bottleneck of
multi-loop calculations.  The standard approach is the Laporta algorithm, which
  systematically solves IBP identities~\cite{Chetyrkin:1981qh,Laporta:2001dd} to express
  any integral in terms of MIs.  Despite decades of improvement
  (e.g.\ Air~\cite{vonManteuffel:2012np}, Kira~\cite{Kira:2021dvi},
  Firefly~\cite{Wallet:2023aa}), the computational cost
remains prohibitive for diagrams beyond two loops.

The Large Index Expansion provides an alternative paradigm.  Instead of directly
solving IBP equations, it exploits the \emph{asymptotic structure} of Feynman integrals
in the limit of large indices.  By expanding $J_{n\theta_a+\nu_a}$ as $n\to\infty$ and
reconstructing the expansion coefficients, one can read off reduction rules that are
valid for \emph{all} indices simultaneously.  This is the approach we develop,
implement, and document here.

% ==============================================================================
%  Section 2: Large Index Expansion Formalism
% ==============================================================================
\section{Large Index Expansion: The Formalism}\label{sec:formalism}

\subsection{The asymptotic ansatz}

Consider a Feynman integral with propagator indices $(n_a)_{a=1}^N$:
\[
  J_{n_1,\ldots,n_N}
  = \int [dl] \frac{1}{D_1^{n_1}\cdots D_N^{n_N}},
  \qquad [dl]=\frac{d^d l}{i\pi^{d/2}} .
\]
For a given sector $\theta_a\in\{0,1\}$ we substitute
\[
  \nu_a \;\to\; \theta_a n + \nu_a ,
\]
and consider the limit $n\to\infty$.  The leading asymptotic behaviour takes the form
\[
  J_{\theta_a n+\nu_a}
  \;\sim\; C(n)\;g_{\underline{\nu}}(n)
  \;=\; C(n)\!\left(\prod_{a=1}^{N}A_a^{\nu_a}\right)
     \sum_{k=0}^{\infty}\frac{h_k(\underline{\nu})}{n^{k}},
  \qquad n\to\infty .
\tag{2.1}\label{eq:ansatz}
\]
Here:
\begin{itemize}
  \item $C(n)$ is a prefactor depending on $n$ and on which propagators carry
        the large index (i.e. on the sector $\theta_a$);
  \item $A_a$ are algebraic numbers (the \textbf{characteristic roots}) determined by
        the sector's IBP system;
  \item $h_k(\underline{\nu})$ is a polynomial in the shifted indices
        $\underline{\nu}=(\nu_1,\ldots,\nu_N)$ of degree $l(k)=rk$
        (typically $r=1,2,3$):
        \[
          h_k(\underline{\nu})
          = \sum_{|\underline{\alpha}|\le l(k)}
            c_{\underline{\alpha}}^{(k)}\,
            \underline{\nu}^{\underline{\alpha}} .
        \tag{2.2}\label{eq:poly}
        \]
\end{itemize}

Equation~\eqref{eq:ansatz} is inserted into the IBP equations written in
momentum representation.  Because the expansion is performed \emph{before} solving the
IBP system, the IBP equations constrain the unknown parameters $(A_a,\,c_{\underline\alpha}^{(k)})$
without ever requiring the explicit form of the integrals.

\subsection{Three-step pipeline}

From the ansatz we obtain a three-stage computational pipeline:

\[
\boxed{
  \text{solve characteristic eq.}\;A_a
  \;\xrightarrow{\quad}\;
  \text{linear eq.}\;c_{\underline\alpha}^{(k)}
  \;\xrightarrow{\quad}\;
  \text{linear eq.}\;b_{\underline\nu}(\underline\nu,n)
}
\tag{2.3}\label{eq:pipeline}
\]

\begin{enumerate}
  \item[\ref{eq:pipeline}a] \textbf{Characteristic equation.}  The leading-order
        ($k=0$) term in~\eqref{eq:ansatz} plugged into the IBP yields algebraic
        equations for $\{A_a\}$, the \textbf{characteristic equation}.
  \item[\ref{eq:pipeline}b] \textbf{Expansion coefficients.}  With $A_a$ known,
        the higher-order coefficients $c_{\underline\alpha}^{(k)}$ are obtained from
        linear equations.  This is the most computationally expensive step and is
        the focus of \S\ref{sec:expansion}.
  \item[\ref{eq:pipeline}c] \textbf{Relation reconstruction.}  A reduction-ansatz
        is assumed for the target integral,
        \[
          \sum_{|\beta|\le m} b_{\underline\nu-\beta}(\underline\nu,n)\,
          g_{\underline\nu-\beta}(\underline\nu,n)=0,
          \qquad \deg_{\nu}b\le k,
        \tag{2.4}\label{eq:ansatz_reduction}
        \]
        and substitution of the known expansion~\eqref{eq:ansatz} yields a linear
        system for the coefficients $b$.
\end{enumerate}

% ==============================================================================
%  Section 3: Solving the Characteristic Equation
% ==============================================================================
\section{Solving the Characteristic Equation}\label{sec:characteristic}

The characteristic equation originates from the $O(n)$ terms in the IBP equations
after inserting~\eqref{eq:ansatz}.  In the original formulation the equations are
rational functions of $\{A_a\}$ and have high algebraic degree because of
denominator cancellation.  Three independent improvements were developed.

\subsection{Inverse-variable substitution}

Introduce $B_a\equiv 1/A_a$.  The $2N$ variables $(A_a,B_a)$ satisfy the bilinear
system
\[
  \sum_{b,a}F_{ba}^{(m)}B_bA_a + F_a^{(m)}A_a = 0,
  \qquad A_aB_a-1=0,
  \qquad \forall a,\; \forall m ,
\tag{3.1}\label{eq:bilinear}
\]
which remains of uniform degree 2 and admits floating-point evaluation before
switching to a finite field for exact arithmetic.

\subsection{Groebner basis with prime decomposition}

Computing a Groebner basis under a lexicographic ordering produces a chain of
ideals; each prime component is then extracted via
$\texttt{minAssGTZ}$ (Singular~\cite{DGPS}).  Only the minimal prime ideals
correspond to distinct solutions, so the expensive computation of embedded
primes is avoided.

\subsection{Sector-limit decomposition}

A key observation is that taking limits in which subsets of propagators are
``switched off'' ($\nu_a\to\infty$ with $\theta_a=0$) reduces the complexity of
the characteristic equation stepwise.  For $\theta_a\in\{0,1\}$ the system
becomes
\[
  \sum_{b,a}F_{ba}^{(m)}B_bA_a + F_a^{(m)}A_a = 0,
  \qquad
  \begin{cases}
    A_aB_a-1=0, & a\in S,\\[2pt]
    A_a=0,      & a\notin S,
  \end{cases}
  \qquad m=1,\ldots,n_{\rm IBP},
\tag{3.2}\label{eq:sector_limit}
\]
where $S\subseteq\{1,\ldots,N\}$ labels the active propagators.  The roots of
simpler subsectors map to a subset of the full sector's roots; the remaining
roots correspond to $A_a\to\infty$.  Different homotopy paths in the space of
solutions give different permutations of the root-to-subsector correspondence.
Because the expansion is needed only at sufficiently high order, any valid
permutation yields a correct reconstruction.

% ==============================================================================
%  Section 4: Expansion Coefficients: Block-Recursive Structure
% ==============================================================================
\section{Expansion Coefficients: Block-Recursive Structure}\label{sec:expansion}

With $\{A_a\}$ known, the coefficients $c_{\underline\alpha}^{(k)}$ of the
polynomials $h_k$ are determined by linear equations derived from the IBP.
A major finding is that these equations possess a \textbf{block-triangular
structure} that can be exploited algorithmically.

\subsection{Recursion from IBP}

Starting from the IBP operator
\[
  \bigl(F_d^{(m)}+F_{ba}^{(m)}1_b^{-}1_a^{+}+F_a^{(m)}1_a^{+}\bigr)g_{\nu}=0,
\tag{4.1}\label{eq:IBP_g}
\]
with $1_a^{+}g_{\nu}=\nu_a g_{\nu+e_a}$, $1_a^{-}g_{\nu}=g_{\nu-e_a}$, and
substituting $g_{\nu}=A^{\underline\nu}h_{\nu}$ with the large-index substitution
$\nu_a\to n\theta_a+\nu_a$, one obtains at order $1/n^{k+1}$ the relation
\[
  \sum_{n\ge1}\sum_{|\beta|=n}(M_{+n}^{(m)}(\alpha))_{\beta}\,
    c_{\alpha+\beta}^{(k+1)}
  = -\!\sum_{b}(N_{-1}^{(m)})_b\,c_{\alpha-e_b}^{(k)}
   -\!\sum_{n\ge0}\sum_{|\beta|=n}(N_{+n}^{(m)}(\alpha))_{\beta}\,
    c_{\alpha+\beta}^{(k)} ,
\tag{4.2}\label{eq:recursion}
\]
where the coefficient matrices are
\[
  (M_{+1}^{(m)})_a
    =\Bigl(\sum_b\tilde F_{ba}^{(m)}+\tilde F_a^{(m)}\Bigr)\theta_a
     -\sum_b\tilde F_{ab}^{(m)}\theta_b,
  \quad
  (N_{-1}^{(m)})_a
    =\sum_b\tilde F_{ba}^{(m)}+\tilde F_a^{(m)},
\]
with $\tilde F_{ba}^{(m)}=F_{ba}^{(m)}A_b^{-1}A_a$,
$\tilde F_a^{(m)}=F_a^{(m)}A_a$.

Equation~\eqref{eq:recursion} is a block-triangular system: for fixed $k$,
coefficients with $|\alpha|=l(k+1)-i$ are solved before those with
$|\alpha|=l(k+1)+1-i$, and the homogeneous part at every step is governed by
$M_{+1}$ alone.

\subsection{Seed-based block solving}\label{ssec:seed_block}

Define the seed indices $\alpha_s$ for which the left-hand side of~\eqref{eq:recursion}
is \emph{empty}; then the right-hand side determines $c^{(k+1)}_{\alpha_s+e_a}$.
For $l(k)=2k$ (the generic case) the seed sequence is:
\[
  |\alpha_s|=2k+1 \;\ra\; c^{(k+1)}_{\alpha_s+e_a}
  \;\ra\; |\alpha_s|=2k \;\ra\; \cdots
\tag{4.3}\label{eq:seed_seq}
\]
The seed indices form the faces of an $N$-dimensional simplex:
\[
  \alpha_s=(l-i-j,i,j,0,\ldots,0),\quad j>0,\quad
  |\alpha_s|=l .
\tag{4.4}\label{eq:simplex}
\]
Geometrically, the solving cascade follows the inclusion chain of faces
$\sigma_0\subset\sigma_1\subset\cdots\subset\sigma_N$ of this simplex.

% ==============================================================================
%  Section 5: Classification of Solution Spaces
% ==============================================================================
\section{Classification of Solution Spaces}\label{sec:classification}

The index-growth rate $r=l(k)/k$ and the dimension $p$ of the space of solutions
for $h_k(\nu)$ are not arbitrary: they are determined by algebraic properties of
the characteristic roots.

\subsection{Rank classification}\label{ssec:rank_class}

The pair of matrices $(M_{+1},N_{-1})$ controls the low-order behaviour.
Three cases occur:
\begin{enumerate}
  \item $\operatorname{corank}M_{+1}=0$, $N_{-1}\equiv0$:\quad $l(k)=k$,
        one or more independent solutions may exist depending on $M_{+1}$;
  \item $\operatorname{corank}M_{+1}=0$, $N_{-1}\not\equiv0$:\quad $l(k)=2k$,
        unique solution (generic case);
  \item $\operatorname{corank}M_{+1}>0$:
        \begin{itemize}
          \item if $(M_{+1})X=-(N_{-1})$ is solvable, $l(k)=2k$, multiple solutions;
          \item otherwise $l(k)\ge 3k$, multiple solutions.
        \end{itemize}
\end{enumerate}
The ratio $r=l(k)/k$ can therefore be $1$, $2$, or $\ge 3$.

\subsection{Geometric meaning}\label{ssec:geometric}

Using $z_a=B_a=1/A_a$, the characteristic equation at a given sector can be
written as
\[
  \sum_{a} g_a^{(m)}(z)\,A_a = 0,
  \qquad
  g_a^{(m)}(z) \equiv \sum_{b} F_{ba}^{(m)}\,z_b + F_a^{(m)},
\tag{5.1}\label{eq:char_z}
\]
where the $F$-matrices are those of the momentum-space IBP system.
The two leading recursion matrices are then obtained by extracting the
coefficients of $\nu_a$ and $1$ from the $O(n)$ terms of the IBP equation after
inserting the ansatz $g_\nu=A^{\underline\nu}h_\nu$ and expanding in $1/n$:
\[
  (N_{-1}^{(m)})_a
    = \left.\frac{g_a^{(m)}(z)}{z_a}\right|_{z=z_{\rm ch}},
\tag{5.2a}\label{eq:N_-1_z}
\]
\[
  (M_{+1}^{(m)})_a
    = \Bigl(\bigl(M_{+1}^{\rm sym}\bigr)^{(m)}\Bigr)_a
      \;\equiv\;
      \Bigl(\Bigl.\sum_b
        \frac{\partial g_b^{(m)}}{\partial z_a}(z)\Bigr|_{z=z_{\rm ch}}\Bigr)
      - \theta_a\!\sum_b\!
        \Bigl(\frac{F_{ba}^{(m)}}{z_b}
          -\delta_{ab}\,\frac{g_b^{(m)}(z)}{z_b^{2}}\Bigr)\Big|_{z=z_{\rm ch}}} .
\tag{5.2b}\label{eq:M_+1_z}
\]

\noindent\textbf{Derivation of \eqref{eq:M_+1_z}.}  From the IBP operator
$P_m = F_d^{(m)}+F_{ba}^{(m)}1_b^{-}1_a^{+}+F_a^{(m)}1_a^{+}$ acting on
$g_\nu=A^{\underline\nu}h_\nu$, the $O(n)$ contribution arises when the raising
operator $1_a^{+}$ hits either the $A^{\underline\nu}$ prefactor (producing the
ratio $A_a/A_b$) or the $h_\nu$ polynomial (producing the factor $\nu_a$).
Extracting the coefficient of $c_{\alpha+e_a}^{(k+1)}$ from the left-hand side
of the IBP gives precisely the expression $(M_{+1}^{\rm sym})_a^{(m)}$ above;
the $-\theta_a$-term accounts for the fact that $\nu_a=n\theta_a+\nu_a^{0}$
contains a piece proportional to $n$ that survives even when the shifted
indices $\nu^{0}$ are held fixed.  Full details are given in the December 2025
monthly report, Sec.~4.

Two geometric consequences follow:
\begin{enumerate}
  \item $\operatorname{corank}M_{+1}>0$ signals that the hypersurfaces
        $\sum_a g_a^{(m)}(z)A_a=0$ intersecting at the characteristic point have
        coalesced --- i.e.\ the characteristic equation acquires multiple roots.
        In this case the expansion space gains independent solutions, one per
        additional sheet of the Borel transform.
  \item $N_{-1}=0$ (i.e.\ $g_a^{(m)}(z_{\rm ch})=0$ for all $a\in\text{PD}$)
        forces $B_a=0$ ($A_a\to\infty$) via the syzygy
        $g_a\partial_aB=g_0B=0$.  This corresponds to collinear external momenta
        or vanishing Gram determinants (second-class singularities).  In these
        cases the degree growth simplifies to $r=1$.
\end{enumerate}

An explicit example is the family $(s,m_1^2,m_2^2)=(s;-1,0)$; at $s=0$ one
finds $p=2$, $r=2$, while at $s=8$ (where the two characteristic roots merge)
$\operatorname{corank}M_{+1}>0$ and $r=3$, confirming that root-coalescence
directly raises the index-growth rate.

\subsection{H${}^{\mathbf{(mn)}}$ rank constraint}\label{ssec:Hmn}

Define
\[
  H^{(mn)}=\sum_a (N_{-1}^{(m)})_a\,(M_{+1}^{(n)})_a,
\tag{5.3}\label{eq:Hmn}
\qquad
  (M_{+1})_a^{(m)}X_{ab}=(N_{-1})_b^{(m)} .
\]
When $M_{+1}$ is full-rank, the non-degenerate subspace of $H$ constrains the
space of possible $c_{\alpha}^{(k)}$ solutions.  The rank of $H$ has been
tabulated for all one- and two-loop families:

\begin{table}[htbp]
  \centering
  \begin{array}{c|c|c|c|c}
    \hline
    \text{family} & n_{\rm prop} & n_{\rm IBP} & \operatorname{rank}M & \operatorname{rank}H \\
    \hline
    \text{1-loop $n$-pt} & n & n & n & n-1 \\
    \text{2-loop vacuum} & 3 & 4 & 3 & 2 \\
    \text{2-loop 2-pt} & 5 & 6 & 5 & 4\;(2)^{\dagger} \\
    \text{2-loop 3-pt} & 7 & 9 & 7 & 6\;(4)^{\dagger} \\
    \text{2-loop 4-pt} & 9 & 10 & 9 & 6 \\
    \text{2-loop 5-pt} & 11 & 12 & 11 & 8 \\
    \hline
  \end{array}
  \caption{Rank of $M_{+1}$ and $H^{(mn)}$ for representative loop families.}
  \label{tab:Hrank}
\end{table}
\noindent $^{\dagger}$Each parenthesised value is the rank of $H$ for a distinct
characteristic root (i.e. a distinct subsector of the same family); the first
value corresponds to the top-sector root.

A change of variables $\tilde\alpha_a=\Lambda_{ab}\alpha_b$ exploiting the
null-space of $H$ typically reduces the number of independent coefficients by
approximately one half.

% ==============================================================================
%  Section 6: Two-Step Reduction Framework
% ==============================================================================
\section{Two-Step Reduction Framework}\label{sec:reduction}

Having obtained the asymptotic expansion, we now reconstruct reduction rules.
The criteria for an \emph{appropriate} set of relations are:
\begin{enumerate}
  \item \textbf{Generativity}: the relations generate all IBP, i.e.\ every integral
        can be reduced to MIs;
  \item \textbf{Solvability}: the relations form a recursive substitution system,
        or can easily be cast into one.
\end{enumerate}
We satisfy both by a \textbf{two-step procedure}.

\subsection{Step 1: Minimal generating set $S$}

Define the \textbf{cone} (reverse pyramid) of height $m$:
\[
  C_m = \{\alpha\in\Z_{\ge0}^N \mid |\alpha|\le m\}.
\tag{6.1}\label{eq:cone}
\]
We search over increasing $m$ and polynomial degree $k$ for a set $S\subset C_m$
such that
\[
  \sum_{\alpha\in C_m} b_{\alpha}(\nu)\,j_{\nu-\alpha}=0,
  \qquad \deg_{\nu}b_{\alpha}=k,
\tag{6.2}\label{eq:search}
\]
generates all IBP identities.  In practice (1-loop and simple 2-loop families)
the minimal values are $m\le2$, $k\le1$.

During the search, solved variables are eliminated so that the leading terms of
the generators do not divide each other --- precisely the Groebner basis
condition in the non-commutative algebra
\[
  \langle S\rangle
  =\Bigl\langle\sum_{\alpha}b_{\alpha}(\nu)
              \prod_a(\Delta_a^{-})^{\alpha_a}\Bigr\rangle
  \subset \mathbb C[\underline\nu,\underline\Delta^{-}] .
\tag{6.3}\label{eq:noncomm_gb}
\]
A useful empirical observation: the generating set found at a given $(m,k)$ is
strictly smaller than the Groebner basis that would be produced by a
brute-force search, i.e.\ $\#S < \#\text{GB}$.

\subsection{Step 2: Block building on a cone}

With $S$ in hand, we build block relations on a larger cone $C_h$:
\[
  0 = \sum_{\beta\in C_{h-m}}\;\sum_{\alpha\in C_m}
      b_{\alpha}^{(i)}(\nu-\beta)\,j_{\nu-\beta-\alpha},
\tag{6.4}\label{eq:block}
\]
where $\nu-\beta$ are called the \textbf{seeds}.  Three block types are needed
for different rank regions:

\begin{enumerate}
  \item[$B_1$ (high rank):] $J_{B_1}=C_{h_1}$; increase $h_1$ until top-rank
        integrals $j_{\nu-\alpha}$, $|\alpha|=h_1$, are reducible.
  \item[$B_2$ (middle rank):] Fix $\underline\nu_{\rm ISP}=0$; increase $h_2$
        until all ranks are reducible for large $\nu_{\rm PD}$.
  \item[$B_3$ (near-corner):] Build on $C_{h_3}$ with
        $\nu_{\rm PD}=1$, $\nu_{\rm ISP}=0$ and add seeds
        $\beta_{\rm super}=-\sum_a\theta_a e_a$ to obtain an extended cone.
\end{enumerate}

A block is \textbf{solvable} if for every target index $\alpha_0$ in a chosen
target set $(J_B)_{\rm targ.}$ there exist relations
\[
  b^{(\alpha_0,i)}(\nu)\,j_{\nu-\alpha_0}
  = \sum_{\nu-\beta<\nu-\alpha_0} b_{\beta}^{(i)}(\nu)\,j_{\nu-\beta},
\tag{6.5}\label{eq:solvable_block}
\]
i.e.\ the integral at the target can be expressed in terms of strictly simpler
integrals.  Finding the minimal set of such relations is equivalent to computing
a Groebner basis of the module
\[
  M_B = \Bigl\langle\sum_{\alpha}b^{(i)}(\nu)_{\alpha}\,e_{\alpha}
                \Bigr\rangle_{i=1}^{|B|}
        \subset \mathbb C[\nu]^{|J_B|}
\tag{6.6}\label{eq:module_gb}
\]
in a position-over-term ordering.

% ==============================================================================
%  Section 7: Shape Unification of IBP Generators
% ==============================================================================
\section{Shape Unification: IBP, Syzygy, and Symbolic Block-Triangular Forms}\label{sec:shapes}

Three apparently distinct families of IBP-generating sets are revealed to be
instances of a single geometric template.

Define the \textbf{extended cone} with respect to a propagator set
$\theta\subset\{1,\ldots,N\}$:
\[
  C_m(\theta) =
  C_m \;\cup\;
  \{-e_a+\alpha \mid a\in\theta,\;\alpha\in\Z_{\ge0}^N,\;|\alpha|\le m+1\} .
\tag{7.1}\label{eq:extended_cone}
\]

All known IBP generator shapes fit into $C_m(\theta)$:
\[
\begin{aligned}
  S(\text{original IBP}) &= C_0(\{1,\ldots,N\}),\\
  S(\text{syzygy IBP in sector }\theta) &= C_m(\text{ISP}_{\theta}),\\
  S(\text{symbolic BT form}) &= C_m(\varnothing).
\end{aligned}
\tag{7.2}\label{eq:shape_family}
\]

Furthermore:
\begin{itemize}
  \item Original IBP is the zero-sector limit of syzygy IBP
        ($\text{PD}=\varnothing$, $\text{ISP}=\{1,\ldots,N\}$, level $=0$);
  \item Symbolic block-triangular (BT) form has the same shape as syzygy IBP in
        the maximal sector ($\text{PD}=\{1,\ldots,N\}$, $\text{ISP}=\varnothing$);
  \item Any symbolic BT form with coefficients linear in $\nu$ can be rewritten in
        syzygy form via normal ordering $1_a^{+}\cdots1_b^{-}$.
\end{itemize}
This unification implies that the cone $C_m(\theta)$ is in some sense the
universal shape underlying all practical IBP reduction methods.

% ==============================================================================
%  Section 8: Completeness of Asymptotic Solutions
% ==============================================================================
\section{Completeness of Asymptotic Solutions}\label{sec:completeness}

A subtle question: does the asymptotic expansion constructed from a single
sector limit generate reduction rules that are valid for \emph{all} sectors?
Two empirical observations, backed by partial proofs, suggest an affirmative
answer.

\begin{enumerate}
  \item[\textbf{C1}] In the sector limit $\nu_a\to\theta_an+\nu_a$, the number of
        solutions to the characteristic equation (counted without multiplicity)
        equals the number of master integrals in that sector and all its
        subsectors.

  \item[\textbf{C2}] Reduction rules reconstructed from a single sector limit
        successfully reduce integrals in every other sector.
\end{enumerate}

\subsection*{Microlocal-analysis justification of C1}

The space of master integrals is linearly isomorphic to a quotient module of the
Weyl algebra:
\[
  M = \mathbb C[z_a]\langle\partial_a\rangle\Big/\!\sim,
\tag{8.1}\label{eq:weyl_quotient}
\]
where the equivalence relation $\sim$ imposes IBP generators
($\hat P_j = F_d^{(j)}-\sum_a(\sum_bF_{ba}^{(j)}z_b+F_a^{(j)})\partial_a$),
ISP constraints ($\partial_a=0$ for $a\in\text{ISPs}$), and the
PD constraints ($\sum_aF_a^{(j)}\partial_a+1$ for $a\in\text{PDs}$).
\textbf{Note:} the shorthand notation $(c_e\partial_e+1)$ used in the main text
suppresses the index structure and $F$-coefficients; the precise form is
$(\sum_aF_a^{(j)}\partial_a+1)$ acting on the propagator derivatives.

The number of independent solutions equals the number of points in the
\textbf{characteristic variety} ---
the algebraic set obtained by replacing $\partial_a$ with commutative variables
$\xi_a$ in the highest-order part of the differential operator:
\[
  \sum_a\!\Bigl(\sum_bF_{ba}^{(j)}z_b+F_a^{(j)}\Bigr)\xi_a=0,\;
  \xi_a=0\;(a\in\text{ISPs}),\;
  z_a\xi_a=0\;(a\in\text{PDs}) .
\tag{8.2}\label{eq:char_variety}
\]

The asymptotic characteristic equation differs only by the homotopy
$z_a\xi_a=0\;\to\;z_a\xi_a=1$ for $a\in\text{PDs}$, which does not change the
number of solutions (generically).  Hence the count of asymptotic solutions
matches the count of master integrals.

% ==============================================================================
%  Section 9: Implementation
% ==============================================================================
\section{C++ Implementation}\label{sec:implementation}

The algorithm was first prototyped in \emph{Mathematica} and then rewritten in
C++ for performance.  Three engineering milestones were reached.

\subsection{Finite-field arithmetic (FFInt)}

To avoid floating-point precision loss in rational reconstruction, all
computations were migrated to the finite field $\mathbb F_p$ with a large prime
$p$ (e.g.\ $p=179\,424\,673$).  The \texttt{firefly::FFInt} type provides
addition, subtraction, multiplication, inversion, and is compatible with the
Eigen linear-algebra backend.  IBP matrix data is loaded from binary files that
store the prime at the header, ensuring field consistency across runs.

\subsection{CMake build system}

A unified \texttt{CMakeLists.txt} replaces hand-rolled build scripts.
Eigen3 and Firefly are auto-detected; the build system generates platform-native
project files and supports out-of-source builds, unit-test execution, and
benchmarking targets.

\subsection{Performance benchmarks}

The C++ implementation of the expansion-coefficient generator
(\texttt{test\_expandFF.cpp}) achieves a $20$--$30\times$ speedup over the
prototype \emph{Mathematica} code.  For the Double Box topology at order $k=4$,
a single sector's single characteristic-root computation completes in a few
seconds.  A coefficient-caching layer (\texttt{SeriesCoefficientIO}) serialises
results to binary for reuse by the relation solver.

The relation solver (\texttt{test\_RelationFF.cpp}) loads cached coefficients,
assembles a global linear system at random sampling points in $\nu$-space,
and calls \texttt{solveLinearSystem} to obtain the $b$-coefficients.
Future work will target sparse-matrix storage to handle more complex topologies.

% ==============================================================================
%  Section 10: Summary and Outlook
% ==============================================================================
\section{Summary and Outlook}\label{sec:summary}

We have presented a complete account of the Large Index Expansion method for
Feynman integral reduction.  The key results are:

\begin{enumerate}
  \item The asymptotic ansatz \eqref{eq:ansatz} converts the IBP system into a
        hierarchical linear problem with three stages
        \eqref{eq:pipeline}: characteristic equation, expansion coefficients,
        and relation reconstruction.
  \item The characteristic equation is tractable via three improvements:
        inverse-variable substitution \eqref{eq:bilinear},
        Groebner-basis prime decomposition, and sector-limit decomposition
        \eqref{eq:sector_limit}.
  \item The expansion coefficients obey a block-recursive structure
        \eqref{eq:recursion} governed by seed indices \eqref{eq:seed_seq}
        forming an $N$-simplex \eqref{eq:simplex}.
  \item The solution-space dimension and growth rate $r=l(k)/k$ are classified
        by the rank of $(M_{+1},N_{-1})$; the $H^{(mn)}$ matrix
        \eqref{eq:Hmn} further constrains the solution space, reducing the
        independent coefficient count by $\sim 50\%$.
  \item The two-step reduction framework (generating set $S$ + solvable blocks $B$)
        on cone shapes $C_m$ satisfies both the generativity and solvability
        criteria.
  \item All IBP generator families (original, syzygy, symbolic BT) are unified
        as extended cones $C_m(\theta)$ \eqref{eq:extended_cone}.
  \item Asymptotic solution completeness follows from the microlocal analysis
        of the characteristic variety \eqref{eq:char_variety}.
  \item The C++/finite-field implementation achieves a $20$--$30\times$
        speedup over the \emph{Mathematica} prototype.
\end{enumerate}

\textbf{Open directions:}
\begin{itemize}
  \item Sparse-matrix optimisation for the relation solver to handle
        three-loop and higher topologies.
  \item Generalisation to integrals with irrational kinematics (roots of
        kinematic invariants) via coordinate-ring extensions
        $\mathbb Q[r_1,\ldots,r_n]$.
  \item Extension to the Large Dimension Expansion ($d\to\infty$) and
        hybrid LIE-LDE strategies.
  \item Systematic comparison with state-of-the-art tools (Kira, Firefly)
        on a common benchmark set.
\end{itemize}

% ==============================================================================
%  References
% ==============================================================================
\bibliographystyle{plain}
\bibliography{refs}

\end{document}
